% Biên dịch: xelatex uct_math_vi.tex
\documentclass[conference]{IEEEtran}
\usepackage{fontspec}
\usepackage[english]{babel}
\babelprovide[import]{vietnamese}
\usepackage{amsmath,amssymb,amsthm,mathtools}
\usepackage{hyperref,algorithm,algpseudocode,enumitem}
\setmainfont{Liberation Serif}
\hypersetup{colorlinks=true,linkcolor=blue,citecolor=blue,urlcolor=blue}
\newtheorem{definition}{Định nghĩa}[section]
\newtheorem{proposition}[definition]{Mệnh đề}
\newtheorem{theorem}[definition]{Định lý}
\newtheorem{remark}[definition]{Nhận xét}
\newcommand{\E}{\mathbb{E}}
\newcommand{\Prob}{\mathbb{P}}
\title{Cơ sở toán học và phân tích thuật toán UCT\\
\large Ứng dụng trong Monte Carlo Tree Search cho Tic-Tac-Toe/Caro}
\author{\IEEEauthorblockN{Tên sinh viên}
\IEEEauthorblockA{Bộ môn Trí tuệ nhân tạo\\
Trường/Đơn vị công tác}}
\begin{document}
\maketitle
\begin{abstract}
Upper Confidence Bound applied to Trees (UCT) là quy tắc chọn nút của Monte
Carlo Tree Search (MCTS), được xây dựng từ bài toán multi-armed bandit. Bài viết
trình bày mô hình trò chơi hữu hạn, suy ra công thức UCB1 bằng cận Hoeffding,
giải thích hằng số khám phá và bốn pha Selection--Expansion--Simulation--
Backpropagation. Bài viết cũng ánh xạ công thức vào chương trình Caro 5 quân
hiện có và đưa ra ứng dụng bonus cho Tic-Tac-Toe 3x3. Không có số liệu thực
nghiệm nào được tự tạo khi chưa chạy thực nghiệm.
\end{abstract}
\begin{IEEEkeywords}
MCTS, UCT, UCB1, multi-armed bandit, trò chơi đối kháng, Tic-Tac-Toe, Caro
\end{IEEEkeywords}

\section{Giới thiệu}
Monte Carlo Tree Search xây dựng dần cây tìm kiếm bằng các lần mô phỏng thay vì
mở rộng toàn bộ cây trò chơi. UCT là quy tắc Selection nổi tiếng: nó ưu tiên nút
có trung bình phần thưởng cao nhưng vẫn dành ngân sách cho nút ít được thăm
\cite{kocsis2006}.

\paragraph{Ghi chú.} Cách gọi ``MCTS kết hợp UCT'' trong đề tài nên hiểu chính
xác là MCTS dùng UCT trong pha Selection. Rollout, điều kiện thắng và quy ước
phần thưởng là các thiết kế độc lập.

\section{Mô hình toán học của trò chơi}
\begin{definition}[Trò chơi đối kháng hữu hạn]
Một trò chơi được mô hình bởi
\[
\mathcal G=(\mathcal S,\mathcal A,T,s_0,\mathcal Z,r),
\]
trong đó $\mathcal S$ là tập trạng thái, $\mathcal A(s)$ là tập nước đi hợp lệ,
$T(s,a)$ là trạng thái kế tiếp, $s_0$ là trạng thái đầu, $\mathcal Z$ là tập
trạng thái kết thúc và $r\in[0,1]$ là phần thưởng nhìn từ một người chơi.
\end{definition}

Với Tic-Tac-Toe, trạng thái gồm bàn cờ 3x3 và người đang đi. Với Caro, trạng
thái có thể là bàn $n\times n$; người chơi thắng khi tạo ít nhất năm quân liên
tiếp theo ngang, dọc hoặc hai hướng chéo. $\mathcal A(s)$ chỉ chứa ô trống;
trạng thái kết thúc khi có người thắng hoặc bàn đầy.

\paragraph{Ghi chú.} Hàm chuyển trạng thái phải bảo toàn luật. Nếu rollout đặt
quân vào ô đã chiếm, các mẫu mô phỏng không còn mô tả đúng trò chơi.

\section{Từ multi-armed bandit đến UCB1}
\subsection{Bài toán bandit}
Xét $K$ cần gạt. Mỗi lần chọn cần $i$ nhận phần thưởng ngẫu nhiên $X_i\in[0,1]$
với kỳ vọng $\mu_i$. Sau $t$ lượt, đặt
\[
N_i(t)=\sum_{\tau=1}^{t}\mathbf 1\{I_\tau=i\},\qquad
\widehat\mu_i(t)=\frac{1}{N_i(t)}\sum_{\tau:I_\tau=i}X_\tau.
\]
Gọi $\mu^*=\max_i\mu_i$ và $\Delta_i=\mu^*-\mu_i$. Regret kỳ vọng là
\[
R_T=\E\!\left[\sum_{t=1}^{T}(\mu^*-\mu_{I_t})\right]
=\sum_{i=1}^{K}\Delta_i\E[N_i(T)].
\]

\subsection{Cận tin cậy và UCB1}
Với biến độc lập trong $[0,1]$, bất đẳng thức Hoeffding cho
\[
\Prob(|\widehat\mu_i-\mu_i|\geq\varepsilon)
\leq 2\exp(-2N_i\varepsilon^2).
\]
Chọn bán kính $b_i(t)=\sqrt{2\ln(t)/N_i(t)}$ dẫn đến chỉ số
\[
U_i(t)=\widehat\mu_i(t)+\sqrt{\frac{2\ln t}{N_i(t)}}.
\]
Thuật toán chọn cần có $U_i(t)$ lớn nhất. Hạng tử đầu là khai thác, hạng tử
thứ hai là khám phá. UCB1 có cận regret logarithmic theo $T$ dưới các giả thiết
bandit chuẩn \cite{auer2002}.

\begin{proposition}[Ý nghĩa UCB]
Khi $N_i(t)$ nhỏ, bán kính tin cậy lớn nên cần $i$ dễ được thử lại. Khi $N_i(t)$
tăng, bán kính giảm theo $N_i(t)^{-1/2}$ và quyết định dần phụ thuộc vào trung
bình thực nghiệm.
\end{proposition}

\paragraph{Ghi chú.} Trong UCT, công thức thường viết
\[
\widehat\mu_j+c\sqrt{\frac{\ln N(s)}{N(s,j)}},
\]
với $c>0$. $c=\sqrt2$ tương ứng một chuẩn hóa phổ biến của UCB1; $c=1.4$ chỉ
là giá trị khởi đầu thực dụng, không phải tối ưu phổ quát.

\section{UCT trên cây trò chơi}
\begin{definition}[Nút tìm kiếm]
Mỗi nút $s$ lưu số lần thăm $N(s)$, tổng phần thưởng $W(s)$, nút cha, các nút
con đã mở rộng và các nước chưa thử. Giá trị thực nghiệm là
\[
Q(s)=W(s)/N(s)\quad\text{khi }N(s)>0.
\]
\end{definition}

Tại nút $s$, với con $s'$ tương ứng nước đi $a$, UCT chọn
\[
a^*\in\arg\max_a\left[Q(s,a)+c\sqrt{\frac{\ln N(s)}{N(s,a)}}\right].
\tag{1}
\]
Nếu $N(s,a)=0$, chỉ số được quy ước bằng $+\infty$ để mọi con mới được thăm.

\subsection{Bốn pha MCTS}
\begin{enumerate}[label=\arabic*.]
\item \textbf{Selection:} đi từ gốc xuống bằng (1) đến nút kết thúc hoặc nút còn
nước chưa thử.
\item \textbf{Expansion:} chọn một nước chưa thử, tạo nút con và áp dụng nước đi.
\item \textbf{Simulation:} rollout đến kết thúc bằng chính sách ngẫu nhiên hoặc
heuristic.
\item \textbf{Backpropagation:} truyền kết quả ngược lên đường đi, cập nhật $N,W$.
\end{enumerate}

\begin{algorithm}[h]
\caption{MCTS--UCT cho trò chơi hai người}
\begin{algorithmic}[1]
\Require trạng thái gốc $s_0$, ngân sách $M$, hằng số $c$
\State tạo nút gốc $v\gets s_0$
\For{$k=1$ to $M$}
  \State $u\gets v$
  \While{$u$ đã mở rộng hết và chưa kết thúc}
    \State $u\gets\arg\max_{x}\left(Q(x)+c\sqrt{\ln N(u)/N(x)}\right)$
  \EndWhile
  \If{$u$ chưa kết thúc và còn nước chưa thử}
    \State tạo một con từ nước chưa thử; đặt $u$ bằng nút mới
  \EndIf
  \State $z\gets\mathrm{Rollout}(u)$; truyền $z$ ngược từ $u$ về $v$
\EndFor
\State \Return con của $v$ có số lượt thăm lớn nhất
\end{algorithmic}
\end{algorithm}

\paragraph{Ghi chú.} Chọn con có nhiều lượt thăm khác chọn con có $Q$ lớn nhất.
Quy tắc max-visits ổn định hơn khi một trung bình cao dựa trên rất ít mẫu; đây
cũng là quy tắc dùng trong mã nguồn hiện có.

\section{Phân tích xác suất và giới hạn}
Tại một nút cố định, nếu coi mỗi hành động là một arm và phần thưởng rollout độc
lập, bị chặn trong $[0,1]$, trực giác UCB1 áp dụng được. Trong cây thật, mẫu được
phân bổ thích nghi và chính sách ở các nút con thay đổi, nên các mẫu không hoàn
toàn độc lập.

\begin{theorem}[Tính nhất quán tiệm cận của UCT]
Trong các giả thiết chuẩn của UCT đối với cây hữu hạn, phần thưởng bị chặn và
rollout hợp lệ, khi số mô phỏng tiến tới vô hạn, xác suất UCT chọn một nước không
tối ưu tại gốc tiến tới $0$ \cite{kocsis2006}.
\end{theorem}

Đây không phải bảo đảm thắng với ngân sách hữu hạn. Cây sâu, phần thưởng hiếm,
nhánh bẫy, rollout thiên lệch và giới hạn thời gian có thể làm kết quả hữu hạn
khác xa giới hạn.

Với $U=Q+c\sqrt{\ln N(s)/N(s,a)}$: $c=0$ thiên về khai thác; $c$ lớn thiên về
khám phá; $c$ trung bình cần hiệu chỉnh thực nghiệm.

\paragraph{Ghi chú.} Không kết luận một $c$ tốt nhất từ một ván. Cần nhiều seed,
đổi người đi trước, ghi tỷ lệ thắng--hòa--thua, thời gian, số nút và đối thủ.

\section{Ánh xạ vào chương trình Caro hiện có}
Trong \texttt{tic\_tac\_toe\_ai.py}, \texttt{MCTSNode} lưu \texttt{visits} và
\texttt{wins}. Hàm \texttt{uct(exploration)} tính
\[
W/N+c\sqrt{\ln N_{\mathrm{parent}}/N},
\]
và trả vô cực cho nút chưa thăm. \texttt{best\_mcts\_move} thực hiện Selection,
Expansion, gọi \texttt{rollout}, rồi Backpropagation theo \texttt{--iterations}.

Phần thưởng nhìn từ người vừa đi là
\[
z=\begin{cases}1,&\text{thắng},\\0.5,&\text{hòa},\\0,&\text{thua}.
\end{cases}
\]
Quy ước phải nhất quán trong UCT và Backpropagation. Rollout heuristic ưu tiên
nước thắng ngay hoặc chặn đối thủ làm thay đổi phân phối mẫu, nên không nên gọi
ước lượng đó là rollout ngẫu nhiên không thiên lệch.

\section{Bonus: Tic-Tac-Toe 3x3}
Với Tic-Tac-Toe, $\mathcal A(s)$ là các ô trống, $T$ đặt quân của người đang đi,
và terminal reward là $1,0,0.5$. Mỗi vòng gồm: mở rộng ô trống; rollout đến khi
có ba quân liên tiếp hoặc bàn đầy; cập nhật đường đi; chọn nước gốc có nhiều lượt
thăm nhất. Tic-Tac-Toe hữu ích vì có thể giải chính xác bằng minimax, do đó UCT
có thể được kiểm tra với đối thủ tối ưu thay vì chỉ đếm một ván thắng.

\section{Thiết kế thực nghiệm đề xuất}
Không đưa số liệu giả. Có thể so sánh $c\in\{0.5,1,\sqrt2,2\}$ và
$M\in\{100,500,1000,5000\}$, dùng ít nhất 30 seed mỗi cấu hình, so sánh rollout
ngẫu nhiên với heuristic, và báo cáo tỷ lệ thắng--hòa--thua, thời gian mỗi nước,
số nút, $N(s_0)$, kích thước bàn, người đi trước, luật thắng và đối thủ.

\section{Kết luận}
UCT đưa cận tin cậy vào cây trò chơi: $Q$ khai thác, còn thành phần căn bậc hai
khám phá. ``Hội tụ với vô hạn mô phỏng'' không đồng nghĩa ``luôn thắng với ngân
sách hữu hạn''. Chất lượng trong Caro/Tic-Tac-Toe phụ thuộc vào sinh nước hợp lệ,
rollout, phần thưởng, $c$ và số vòng lặp.

\begin{thebibliography}{9}
\bibitem{auer2002} P. Auer, N. Cesa-Bianchi, and P. Fischer, ``Finite-time
Analysis of the Multiarmed Bandit Problem,'' \emph{Machine Learning}, 47,
235--256, 2002. \href{https://doi.org/10.1023/A:1013689704352}{doi}.
\bibitem{kocsis2006} L. Kocsis and C. Szepesvári, ``Bandit Based Monte-Carlo
Planning,'' ECML, LNCS 4212, 282--293, 2006.
\href{https://doi.org/10.1007/11871842_29}{doi}.
\bibitem{browne2012} C. B. Browne et al., ``A Survey of Monte Carlo Tree Search
Methods,'' \emph{IEEE TCIAIG}, 4(1), 1--43, 2012.
\href{https://doi.org/10.1109/TCIAIG.2012.2186810}{doi}.
\bibitem{sutton2018} R. S. Sutton and A. G. Barto, \emph{Reinforcement Learning:
An Introduction}, 2nd ed., MIT Press, 2018.
\href{http://incompleteideas.net/book/the-book-2nd.html}{online edition}.
\end{thebibliography}
\end{document}
